% =============================================================================
% Project : THE PHILOSOPHY OF VIABILITY
% Book    : Book I — Ontology of Viability
% Book ID : VS.B01
% Part    : Part I — Identity and Manifestation
% Part ID : VS.B01.P01
% Chapter : Chapter 08 — Comparative Philosophical Examination
% Chapter ID : VS.B01.P01.C08
%
% File :
%   4_ASSETS/tables/
%   VS.B01.P01.C08.S12.SS16.P003.TL01.al_farabi_claims_and_comparative_limits.tex
%
% Section : Section 12 — Al-Farabi
%
% Asset :
%   Table TL01 — Al-Farabi's Claims and Comparative Limits
%
% Structural Identifier :
%   VS.B01.P01.C08.S12.SS16.P003.TL01
%
% Purpose :
%   To distinguish the historical reconstruction of Al-Farabi's
%   principal metaphysical claims from their logical form,
%   epistemic status within the present inquiry, and comparative
%   limits for the Ontology of Viability.
%
% Intellectual Benefit :
%   Prevents the conflation of historical doctrine, interpretive
%   reconstruction, and the foundational commitments of Book I.
%
% Aesthetic / Experiential Benefit :
%   Provides a compact reference for the extended comparative
%   analysis in Section 12.
%
% Representational Limit :
%   The classifications concern the present reconstruction.
%   They are not Al-Farabi's own taxonomy of claims or knowledge.
%
% Dependencies :
%   longtable, booktabs, array.
%
% =============================================================================

\small
\renewcommand{\arraystretch}{1.20}
\setlength{\tabcolsep}{4pt}

\begin{longtable}{
    >{\raggedright\arraybackslash}p{2.0cm}
    >{\raggedright\arraybackslash}p{2.6cm}
    >{\raggedright\arraybackslash}p{2.25cm}
    >{\raggedright\arraybackslash}p{2.55cm}
    >{\raggedright\arraybackslash}p{4.0cm}
}

\caption{Al-Farabi's Principal Claims and Their Comparative Limits}
\label{tab:VS.B01.P01.C08.S12.SS16.P003.TL01} \\

\toprule
\textbf{Claim or issue}
&
\textbf{Historical reconstruction}
&
\textbf{Functional or logical form}
&
\textbf{Epistemic or justificatory status}
&
\textbf{Comparative limit for Book I}
\\
\midrule
\endfirsthead

\toprule
\textbf{Claim or issue}
&
\textbf{Historical reconstruction}
&
\textbf{Functional or logical form}
&
\textbf{Epistemic or justificatory status}
&
\textbf{Comparative limit for Book I}
\\
\midrule
\endhead

\midrule
\multicolumn{5}{r}{\footnotesize\textit{Continued on next page}}\\
\endfoot

\bottomrule
\endlastfoot

Being
&
Al-Farabi distinguishes several senses of Being, including Being as true and Being as circumscribed by a quiddity outside the soul.
&
Proposition
&
Argued or Derived Claims
&
The differentiation of senses does not establish that Being is identical with Identity or that its ontological status depends upon cognition.
\\

\addlinespace

Unity
&
Al-Farabi distinguishes numerical unity, continuity, indivisibility, uniqueness, and circumscription by quiddity.
&
Proposition
&
Argued or Derived Claims
&
The plurality of senses of Unity supports conceptual differentiation but does not establish that Unity is identical with Identity, Singularity, or Composition.
\\

\addlinespace

Unity and quiddity
&
A particular sense of being one is connected with the circumscription of an existent by its quiddity.
&
Proposition
&
Conditionally Supported Claims
&
A specified relation between Unity and quiddity cannot be generalized into an identity of Unity, quiddity, and Identity.
\\

\addlinespace

The First Existent
&
The First is presented as uncaused, independent, simple, unique, and the first cause of the existence of other beings.
&
Foundational Commitment within Al-Farabi's system
&
Conditionally Supported Claims
&
The First Existent is not equivalent to Being as such, Identity, or any proposed realm of Book I.
\\

\addlinespace

Emanation
&
Dependent existents proceed from the First without adding to its perfection or requiring an alteration in it.
&
Thesis
&
Conditionally Supported Claims
&
Causal production is not equivalent to Expression or Manifestation, and causal dependence does not by itself explain the Identity of the dependent existent.
\\

\addlinespace

Secondary intellects and celestial order
&
The First is followed by an ordered series of separate intellects associated with the celestial structure; details differ between relevant treatises.
&
Proposition
&
Conditionally Supported Claims
&
The cosmological hierarchy is not equivalent to the proposed domains and hubs of Book I.
\\

\addlinespace

Active Intellect
&
The Active Intellect occupies a distinct position in the separate-intellect order and is connected with the sublunary domain and human intellectual activity.
&
Proposition
&
Argued or Derived Claims
&
The Active Intellect is not identical with Light or with a foundational component of the ontology of Book I.
\\

\addlinespace

The virtuous city
&
The city is related analogically to the ordered cosmos and to the differentiated organization of an organism.
&
Thesis
&
Conditionally Supported Claims
&
The analogy does not by itself establish the normative legitimacy of a political hierarchy or define the Identity of a social entity.
\\

\addlinespace

Relation to Book I
&
Al-Farabi's analysis of Being, Unity, quiddity, and causal dependence overlaps with problems examined in Part I but uses a distinct explanatory architecture.
&
Thesis
&
Argued or Derived Claims
&
The relation is critical proximity without conceptual equivalence or direct doctrinal continuity.
\\

\addlinespace

Identity and quiddity
&
The relation between Al-Farabi's account of circumscription by quiddity and the ontological status of Identity remains open to further analysis.
&
Open Problem
&
Unresolved or Open Problems
&
The question cannot be settled by substituting Identity for quiddity or by treating the two terms as interchangeable.
\\

\end{longtable}

\normalsize

% =============================================================================
% END OF TABLE TL01
% =============================================================================